\begingroup\small
\begin{longtable}{@{}>{\raggedright\arraybackslash}p{0.06\linewidth}>{\raggedright\arraybackslash}p{0.39\linewidth}>{\raggedright\arraybackslash}p{0.49\linewidth}@{}}
\toprule No. & Family and named implementations & Recurring behavior and proof obligation\\\midrule\endhead
1 & \textbf{Original source / allocation}. Primary PA194 (67/7); PA194 (68/8) & Same original command protocol and separate diagonal calculus. Complete source simulation and PA calibration.\\[5pt]
2 & \textbf{Constructor coefficient}. M4R; M5; M6 & Coefficients four, five and six change numerical formula coding. Hereditary proof recoding, including the zero alias.\\[5pt]
3 & \textbf{Cached MP}. PA194C; PA194CS & Successful MP stores the consequent in cache and clears visible top; Shift tests cache. Adjusted all-store command compiler.\\[5pt]
4 & \textbf{Visible no-push MP}. PA194N; PA194N5; PA194N6 & Successful MP retains (b,w,b,0), omitting dummy-zero push; original sum guard. Reserved-zero handling and arbitrary carried stores.\\[5pt]
5 & \textbf{Cache-only guard}. PA194H; PA194H5; PA194H6; PA194HS = filed PA194NS & Visible no-push MP with Shift test h=1; HS uses symmetric comparison. Final contradiction detection must use Pop.\\[5pt]
6 & \textbf{Symmetric sum guard}. PA194NS (sum guard); PA194SUMS & Original sum condition t+h=1 with symmetric destructive comparison. Source tail equivalence; NS and SUMS differ in MP push behavior.\\[5pt]
7 & \textbf{Q distribution}. PA194Q & Exchange the curried premises of the universal-distribution head. Restricted-rho derivation of premise exchange.\\[5pt]
8 & \textbf{Q with coefficient changes}. PA194Q5; PA194Q6 & Q head plus coefficient five or six; original dummy-push MP. Both premise exchange and radix recoding.\\[5pt]
9 & \textbf{Q with no-push MP}. PA194QN = PA194NQ; PA194QN5; PA194QN6 & Q head, visible no-push MP, original sum guard. Composition of Q proof transport and no-push compiler.\\[5pt]
10 & \textbf{Q with cache-only guard}. PA194QH; PA194QH5; PA194QH6; PA194QHS & Q head, no-push MP, cache-only Shift. Cache-only Pop compiler and Q transport.\\[5pt]
11 & \textbf{Q with symmetric sum guard}. PA194QNS; PA194QS & Q head with the original sum guard; QNS removes MP push, QS retains it. Exact sum-guard block equivalence.\\[5pt]
12 & \textbf{Diagonal choice}. PA194D; PA194DHS; Gamma(D) + CSZA & Use Gamma(D) instead of Gamma(A); accompanying controller variants. New diagonal selector; supplement has narrower evidence.\\[5pt]
13 & \textbf{Known-zero source returns}. PA194A; PA194F; PA194J; PA194ZR; PA194R; PA194CSZ = PA194X; PA194CZ; PA194HA; PA194HAZ; CSZA final-register 3; CSZA final-register 5; CSZA final-register 6 & Replace selected jumps or returns by decrements of registers proved zero. A source invariant and separately verified blank startup are both needed.\\[5pt]
14 & \textbf{Physical register placement}. PA194-AD; PA194F\_AD; PA194R\_AD; PA194TDY; PA194P5; PA194HX; PA194AX11HR & Preserve logical commands while exchanging register positions and routers. Register-frame conjugacy and complete initialization.\\[5pt]
15 & \textbf{Decoder topology}. PA194B; PA1944L & Preserve the original source program. All-tape or complete source-frame positive simulation.\\[5pt]
16 & \textbf{AX11 interface}. PA194AX11 & Push the guard and fall through the full MP operand reader: four operands, eleven entries. All eleven selectors and the fourth operand must be compiled.\\[5pt]
17 & \textbf{AX11 with cache-only guard}. PA194AX11H; PA194AX11QH & Four-operand AX11 with H guard, with original or Q distribution. Final Pop compiler; the revised AX11H manuscript claims the full written iPRA calibration, while its end-to-end Lean formalization remains incomplete.\\[5pt]
18 & \textbf{AX11 with changed coding}. PA194AX11Q6; PA194AX11QH5 & AX11 and Q with coefficient six/sum guard or five/cache guard. Four-field coding, radix transport, appropriate halt detector.\\[5pt]
19 & \textbf{AX11 with symmetric sum guard}. PA194AX11S & Exactly the AX11 sum-guard command semantics. Relative source/physical equivalence.\\[5pt]
20 & \textbf{Explicit generalization binder}. PA194G; PA194AX11G & The first AX operand c controls U(c,h), instead of inherited offset d; ten- or eleven-entry AX. Shorter generalization word with no offset-setting command.\\[5pt]
21 & \textbf{Explicit binder with H guard}. PA194GH; PA194GHS & Explicit binder and ten-entry AX with cache-only Shift. New all-store compiler; complete physical frame checks.\\[5pt]
22 & \textbf{AX11QH decoder alternatives}. PA194AX11QH4L; PA194AX11QH4R & Exactly the existing four-operand AX11QH source. Positive finite-time transfer to the parent.\\[5pt]
23 & \textbf{AX11 with explicit binder and cache-only guard}. PA194AX11GH; PA194AX11GHS & Four-operand eleven-entry AX with U(c,h), Q distribution, no-push MP and cache-only Shift; GHS reverses comparison operands. Written relative iPRA theorem plus inherited revised AX11H absolute calibration; no completed end-to-end Lean theorem.\\[5pt]
24 & \textbf{Direct history selectors}. PA194SEL & Eight-row reorganization of two recurring decoder pairs into history selectors. Direct relative iPRA proof: a parent halt by t gives a SEL halt by 3t, and a SEL halt by t gives a parent halt by t. Absolute calibration is explicitly left open in this packet.\\[5pt]
25 & \textbf{Exact physical three-to-one shortcuts}. PA194AX11GHF; PA194AX11GHSF & Changes native rows 26/1 and 101/0, each matching three parent steps under an incoming-edge invariant. Direct relative iPRA theorem; absolute result uses the same inherited calibration.\\[5pt]
26 & \textbf{Loop-continuation reassociation}. PA194LC; PA194LC\_A; PA194LC\_B; PA194LC\_C; PA194LC\_AX11H & Alternative recurring constructor-loop associations preserve the chosen source graph: 3,673 primary instructions or 3,671 AX11H instructions. Explicit relative iPRA proof to the corresponding parent; the absolute conclusion imports that parent calibration.\\[5pt]
\bottomrule\end{longtable}
\endgroup
